\documentclass{article}
\usepackage{PRIMEarxiv} % Core package for the PrimeAI Template
\usepackage{amsmath, amssymb, amsthm, bm, dcolumn} % Add amsthm here for the proof environment
\usepackage[numbers,sort&compress]{natbib} % Natbib for citations
\usepackage{graphicx} % For high-quality images
\usepackage{multicol} % Multi-column figures/tables
\usepackage{hyperref} % Hyperlinks
\usepackage{listings} % Code listings
\usepackage{authblk} % For structured affiliations
% Define theorem style (optional)
\theoremstyle{plain}
\newtheorem{theorem}{Theorem}
\renewcommand{\qedsymbol}{}

% Custom commands for primes and qubit encoding
\newcommand{\primeQubit}[1]{|p_{#1}\rangle}
\newcommand{\primeGate}[1]{U_{p_{#1}}}

\usepackage{wrapfig}
\usepackage[pscoord]{eso-pic}
\usepackage[fulladjust]{marginnote}
\reversemarginpar

% Typesetting improvements without footnote patching
\usepackage[protrusion=true, expansion=true, tracking=false]{microtype}
\microtypecontext{spacing=nonfrench}

% Line numbers
\usepackage[right]{lineno}

% Text layout - adjust as needed
\raggedright
\setlength{\parindent}{0.5cm}
\textwidth 5.25in 
\textheight 8.75in

% Set double spacing
\usepackage{setspace} 
\doublespacing

% Adjust width for specific content
\usepackage{changepage}

% Adjust caption style
\usepackage[aboveskip=1pt,labelfont=bf,labelsep=period,singlelinecheck=off]{caption}

% Remove brackets from references
\makeatletter
\renewcommand{\@biblabel}[1]{\quad#1.}
\makeatother

% Header, footer, and page numbers
\usepackage{lastpage,fancyhdr}
\pagestyle{fancy}
\fancyhf{}
\fancyhead[L]{Preprint - PrimeAI Enhanced Template}
\fancyfoot[C]{\scriptsize Multiplicity Theory © 2024 Ryan Van Gelder - Citizen Gardens \\ Licensed Under MIT and CC BY-NC-SA 4.0.}
\fancyfoot[R]{Page \thepage\ of \pageref{LastPage}}
\renewcommand{\footrule}{\hrule height 2pt \vspace{2mm}}

\begin{document}
\title{Quantum Aesthetic Theory: Mapping Beauty in Quantum Systems}

\author{Ryan Van Gelder}
\affil{Citizen Gardens - The Foundation of Multiplicity}
\date{\today}
\maketitle

\begin{abstract}
This paper introduces a Quantum Aesthetic Theory (QAT), blending principles of quantum mechanics, multiplicity, and prime-based encoding to objectively quantify and express the harmony of quantum systems. By defining mathematical measures such as symmetry, coherence, and entropy, QAT provides a framework for understanding aesthetic elegance across physics, art, music, and architecture. Applications include quantum-inspired creative expressions and computational models of beauty.
\end{abstract}

\section{Introduction}
The intersection of quantum mechanics and aesthetics presents an intriguing frontier for scientific and artistic exploration. Quantum Aesthetic Theory (QAT) seeks to formalize the notion of "beauty" or "harmony" within quantum systems using mathematical metrics derived from quantum principles and multiplicity theory. This paper explores how aesthetic elegance in quantum systems can be quantified and applied across creative domains such as art, music, and architecture.

\section{Defining Elegance in Quantum Systems}

\subsection{Symmetry and Coherence}
Symmetry and coherence are fundamental indicators of elegance in quantum systems:
\begin{itemize}
    \item \textbf{Symmetry}: Conservation laws and invariants (e.g., angular momentum, parity) provide mathematical representations of beauty in quantum states. Symmetric systems minimize complexity while maintaining order.
    \item \textbf{Coherence}: Defined as the degree of superposition or entanglement, coherence captures the harmonious interplay of quantum states.
\end{itemize}

We define an aesthetic metric \(\aestheticMetric(t)\) for a system as a function of these properties:
\begin{equation}
\aestheticMetric(t) = \coherenceMetric(t) + \alpha \cdot \entropy(t),
\end{equation}
where \(\coherenceMetric(t)\) measures quantum coherence, \(\entropy(t)\) represents disorder, and \(\alpha\) is a scaling factor.

\subsection{Entropy and Order}
Aesthetic harmony often balances order and entropy. Low-entropy quantum systems with high coherence exhibit higher perceived elegance. Mathematically, the relationship can be expressed using the von Neumann entropy:
\begin{equation}
\entropy = - \text{Tr}(\rho \log \rho),
\end{equation}
where \(\rho\) is the density matrix of the quantum system.

\section{Applications Across Domains}

\subsection{Art}
Quantum-inspired art leverages visual representations of quantum coherence and entanglement. Using tensor networks to map quantum states into visual patterns allows for the depiction of dynamic relationships and emergent complexity.

\subsection{Music}
Prime-based encoding can map musical frequencies to quantum states, creating harmonic structures derived from quantum coherence:
\begin{equation}
\primeEncode{i} = 2 \pi \cdot f_{i},
\end{equation}
where \(f_{i}\) represents the frequency of the \(i\)-th note.

\subsection{Architecture}
Architectural designs can be inspired by quantum geometries, using symmetry and fractal-like tensor networks to create scalable and dynamic frameworks. The tensor representation of spatial relationships can be modeled as:
\begin{equation}
T_{ijk} = \sum_{m} \phi_{im} \phi_{jm} \phi_{km},
\end{equation}
where \(T_{ijk}\) encodes multi-scale interactions.

\section{Mathematical Framework}

\subsection{Dynamic Multiplicity Equation}
The time evolution of the aesthetic metric can be described using a dynamic equation:
\begin{equation}
\frac{\partial \aestheticMetric(t)}{\partial t} = \sum_{i=1}^N \alpha_i \aestheticMetric_i(t) + \beta_i \cos(\omega_i t + \phi_i),
\end{equation}
where \(\alpha_i\) and \(\beta_i\) are weighting factors, and \(\omega_i\) and \(\phi_i\) capture phase dynamics.

\subsection{Feedback Mechanisms}
Recursive feedback allows for continuous adaptation of the aesthetic metric. Let \(F(t)\) represent a feedback function:
\begin{equation}
\aestheticMetric(t + \Delta t) = F\big(\aestheticMetric(t), \partial_t \aestheticMetric(t)\big).
\end{equation}

\section{Conclusion}
Quantum Aesthetic Theory establishes a mathematical foundation for exploring beauty in quantum systems and creative domains. By leveraging coherence, symmetry, and feedback, QAT bridges physics with art, music, and architecture. Future work will focus on computational models and real-world implementations to validate these concepts.

\nocite{*}
\bibliographystyle{plain}
\bibliography{references}

\end{document}